\begin{textsample}{POS Dim 2 – ap – Score 26.00 – ap/ap\_em\_35.txt}  \label{ex:f2_pos_088}
EDUCAÇÃO PROFISSIONAL E TECNOLÓGICA INTRODUÇÃO A Educação Profissional e Tecnológica ( EPT ) tem o papel de habilitar os sujeitos \textbf{técnica} e ideologicamente para o trabalho, subordinando as práticas educativas aos interesses do capital, que aponta para as novas demandas de \textbf{qualificação} humana.
E assim, ações governamentais são projetadas e executadas para tal público, dentre as quais, em âmbito \textbf{federal}, gestada por meio de programas e projetos, e realizada por meio da adesão dos estados e \textbf{municípios} ou pela \textbf{instituição} de parcerias público-privado, por sua vez induzidas por pressões diretas do \textbf{governo} \textbf{federal} ou pelas vantagens que proporcionam, especialmente pelos recursos disponibilizados.
De acordo com o Plano Curricular da Educação Básica do Estado do Amapá o \textbf{currículo} do Ensino Médio se organiza em duas partes : a parte de Formação Geral Básica e os \textbf{Itinerários} formativos.
A Formação Geral Básica é composta pelas quatro áreas do conhecimento.
Cada área do conhecimento abrange os saberes específicos de seus componentes curriculares.
Por sua vez, os \textbf{Itinerários} Formativos, além das quatro áreas da Formação Geral Básica, contemplam um quinto \textbf{itinerário} : a Educação Profissional e \textbf{Técnica}.
No Amapá, no âmbito da Secretaria de Estado da Educação ( SEED ), a trajetória da Educação Profissional e Tecnológica ( EPT ) está sendo \textbf{ofertada} em \textbf{instituições} de ensino, Escolas e Centros de Educação Profissional, considerando \textbf{atender} às novas configurações do mundo do trabalho, mas, igualmente, a contribuir para a \textbf{elevação} da \textbf{escolaridade} dos alunos para que possam enfrentar a transformação do mundo contemporâneo.
A nova LDB ( Lei n.o 9.394/1996 ) apesar de apresentar um \textbf{capítulo} específico para a Educação Profissional, seus \textbf{artigos} não cumpriam as finalidades de uma Lei de \textbf{Diretrizes} e Bases, sendo que apenas apontavam \textbf{diretrizes} gerais.
Desta forma, a \textbf{legislação} para a educação profissional foi organizada gradativamente a partir do Decreto no 2.208/97, pelo \textbf{Parecer} CNE/CEB no 16/99 e Resolução no 04/99 CNE/CEB.
Assim a educação profissional foi dividida em três níveis : a ) básico ( \textbf{destinado} à \textbf{qualificação}, requalificação e reprofissionalização dos \textbf{trabalhadores}, independente de \textbf{escolaridade} prévia ) ; b ) \textbf{técnico} ( \textbf{destinado} a proporcionar \textbf{habilitação} profissional a alunos matriculados ou egressos do ensino médio ) ; e, c ) tecnológico ( correspondente a \textbf{cursos} de nível superior na área de tecnológica ), \textbf{destinados} a egressos do ensino médio e \textbf{técnico} ( BRASIL, 1997 ).
A educação tecnológica básica é uma das \textbf{diretrizes} que a LDB estabelece para orientar o \textbf{currículo} do Ensino Médio.
A lei ainda associa a “ compreensão dos fundamentos científicos dos processos produtivos ” ao relacionamento entre teoria e prática em cada disciplina do \textbf{currículo}.
Nessa perspectiva, entende -se que a formação profissional estará contribuindo para a capacitação de pessoas para o mercado, mais especificamente, para a formação do \textbf{trabalhador}, e a escola se apresenta enquanto produtora de capital humano.
Nesta lógica, subordina -se a função social da educação, que passa a ser empreendida com uma função mercadológica, em que seu papel é produzir bons produtos para o sistema econômico.
A implantação do V \textbf{itinerário} formativo requer da educação profissional e tecnológica um modelo de formação amplo, expondo flexibilidade no \textbf{currículo} que permitirá aos estudantes aprofundar os conhecimentos em uma ou mais áreas de seu próprio interesse.
Outro ponto importante é a visão geral dos tipos de \textbf{itinerários}.
Eles estão baseados em áreas do conhecimento ( linguagens e suas tecnologias, matemática e suas tecnologias, ciências da natureza e suas tecnologias, ciências humanas e sociais aplicadas ), ou na formação \textbf{técnica} e profissional.
Já o \textbf{itinerário} formativo integrado combina mais de uma área, podendo incluir a formação profissional.
Por isso, a função do \textbf{itinerário} formativo é contribuir para que o estudante possa traçar rotas e planejar percursos, de forma a conciliar as demandas do mundo do trabalho, as expectativas pessoais e os tipos de formações disponíveis, por isso, a SEED vem mobilizando os recursos indispensáveis para obtenção de resultados satisfatórios, no âmbito da EPT.

% matched lemmas: artigo, atender, capítulo, currículo, curso, destinar, diretriz, elevação, escolaridade, federal, governo, habilitação, instituição, itinerário, legislação, município, ofertar, parecer, qualificação, trabalhador, técnico
\end{textsample}
